% !PRINTED none
% Chapter-opening page: no running head/folio visible in any of the three
% bands (consistent with an Abschnitt-opening page).

\begin{center}
\gesperrt{Sechster Abschnitt.}

\gesperrt{Gruppen linearer Substitutionen.}
\end{center}

\begin{center}
\rule{3cm}{0.4pt}
\end{center}

\begin{center}
§. 37.

\gesperrt{Lineare Substitutionen und ihre Zusammensetzung.}
\end{center}

Eines der wirksamsten Mittel zur Bildung von Gruppen, auf welches zugleich viele Anwendungen führen, sind die linearen Substitutionen und ihre Zusammensetzung. Wir sind schon mehrfach solchen linearen Substitutionen begegnet und haben sie z. B. im zweiten Abschnitte des ersten Bandes bei Gelegenheit des Multiplicationsgesetzes der Determinanten, und sodann im elften Abschnitte bei den äquivalenten Zahlen betrachtet.

Unter einer \gesperrt{linearen Substitution} von $n$ Variablen verstehen wir ein System von Gleichungen, durch das ein System von $n$ Veränderlichen $y_1, y_2, \dots y_n$ linear durch ein anderes System $x_1, x_2, \dots x_n$ ausgedrückt wird. Wir unterscheiden nach der Anzahl der Variablen \gesperrt{unäre, binäre, ternäre, quaternäre} Substitutionen, und wollen im Allgemeinen die Zahl der Variablen die \gesperrt{Dimension} der Substitution nennen. Wir beschränken uns fürs erste auf \gesperrt{homogene Substitutionen}, so dass, wenn mit $a_i^{(\varkappa)}$ die Coëfficienten bezeichnet werden, die Substitution durch das Gleichungssystem

\weq{1}{y_\varkappa = \sum_{1,n}^{i} a_i^{(\varkappa)} x_i \qquad \varkappa = 1, 2, \dots n}

dargestellt ist. Oft kommt es auf die Variablen selbst nicht an, so dass eine solche Substitution durch ihre Coëfficienten $a_i^{(\varkappa)}$ hinlänglich gekennzeichnet ist. Man benutzt zuweilen auch zur Be-
%CONT>>
